\documentclass[10 pt]{article}
\usepackage[utf8]{inputenc}
\usepackage{parskip}
\usepackage[utf8]{inputenc}
\usepackage[spanish]{babel}
\usepackage[left=2cm, right=2cm, top=1.5cm, bottom=2cm]{geometry}
\usepackage{float}
\usepackage{graphicx}
\usepackage{caption}
\usepackage{cancel}
\usepackage{amsmath,amsthm,amsfonts,amscd,amssymb,amsbsy}

\title{\textbf{\underline{Introducción}}}
\author{}
\date{}

\begin{document}

\pagestyle{empty}
\maketitle
\thispagestyle{empty}

\section*{}

\Large{La \textbf{teoría de grupos} nace del estudio de las simetrías y transformaciones que preservan la estructura de los objetos matemáticos. Con el tiempo, los grupos han adquirido un papel central en el álgebra, no solo porque constituyen el punto de partida para el estudio de estructuras más complejas, como los \textbf{anillos}, los \textbf{cuerpos} y los \textbf{módulos}, sino también porque ofrecen un lenguaje unificador para describir fenómenos presentes en numerosas áreas de las matemáticas.} \\

\Large{La noción de \textbf{operación binaria} constituye el punto de partida de la teoría de grupos. En esencia, una operación binaria describe una regla que asigna a cada par de elementos de un conjunto otro elemento del mismo conjunto. La exigencia de que el resultado permanezca dentro del conjunto, expresada por la propiedad de\textbf{ clausura}, permite estudiar la estructura sin abandonar el universo de objetos considerado y será una característica común a todas las estructuras algebraicas que estudiaremos más adelante.} \\

\begin{itemize}
    \item \Large{\textbf{Definición:} Sean $A$ y $G$ conjuntos tales que $G \subseteq A$, y sea $*:A \times A \rightarrow A$ una función. Diremos que $*$ es una \textbf{operación binaria} en $G$, si y sólo si, para todo $a \in G$ y para todo $b \in G$, $*(a,b) \in G$.} \\
\end{itemize}

\Large{\textbf{Nota:} Cuando digamos simplemente que $*$ es una operación binaria en $G$, asumiremos que $*$ es una función de $G\times G$ en $G$.} \\

\begin{flushright}
   $\blacksquare$
\end{flushright}

\Large{\textbf{Ejemplos:}} \\

\begin{itemize}
\item \Large{La suma de números enteros $+:\mathbb Z \times \mathbb Z \rightarrow \mathbb Z$ es una operación binaria en $\mathbb N$.} \\

\item \Large{La resta de números enteros $-:\mathbb Z \times \mathbb Z \rightarrow \mathbb Z$ no es una operación binaria en $\mathbb N$.} \\
\end{itemize}

\begin{flushright}
$\blacksquare$
\end{flushright}

\begin{itemize}
    \item \Large{\textbf{Definición:} Sean $G$ un conjunto y $*$ una operación binaria en $G$. Para todo $a \in G$ y para todo $b \in G$:}
    \begin{align*}
        *(a,b)&=a*b
    \end{align*}
\end{itemize}
\begin{flushright}
    $\blacksquare$
\end{flushright}

\Large{Sean $G$ un conjunto y $*$ una operación binaria en $G$. En vista de la definición anterior, tenemos que para todo $a \in G$ y para todo $b \in G$, $a*b \in G$. A continuación definimos algunas propiedades que puede tener una operación binaria:} \\

\begin{itemize}
    \item \Large{\textbf{Definición:} Sean $G$ un conjunto y $*$ una operación binaria en $G$. Diremos que $*$ es \textbf{asociativa} en $G$, si y sólo si, para todo $a \in G$, para todo $b \in G$ y para todo $c \in G$:}
    \begin{align*}
        a*(b*c)&=(a*b)*c
    \end{align*}
\end{itemize}
\begin{flushright}
    $\blacksquare$
\end{flushright}

\begin{itemize}
    \item \Large{\textbf{Definición:} Sean $G$ un conjunto y $*$ una operación binaria en $G$. Diremos que $*$ es \textbf{conmutativa} en $G$, si y sólo si, para todo $a \in G$ y para todo $b \in G$:}
    \begin{align*}
        a*b&=b*a
    \end{align*}
\end{itemize}
\begin{flushright}
    $\blacksquare$
\end{flushright}

\Large{\textbf{Ejemplo:}} \\

\Large{La composición de funciones es una operación binaria asociativa pero no conmutativa.} \\

\begin{flushright}
$\blacksquare$
\end{flushright}

\Large{Una operación binaria no solo permite combinar elementos de un conjunto, sino que también puede dar lugar a la aparición de \textbf{elementos distinguidos}, caracterizados por las propiedades que satisfacen con respecto a dicha operación. Entre ellos destacan el \textbf{elemento neutro}, que no altera a los demás elementos al operar con ellos, y los \textbf{elementos inversos}, que permiten revertir el efecto de la operación.} \\

\begin{itemize}
    \item \Large{\textbf{Definición:} Sean $G$ un conjunto y $*$ una operación binaria en $G$. Diremos que $G$ tiene \textbf{elemento neutro} respecto a $*$, si y sólo si, existe $e \in G$ tal que para todo $g \in G$, $e*g=g*e=g$.} \\
\end{itemize}

\begin{flushright}
    $\blacksquare$
\end{flushright}

\Large{Sean $G$ un conjunto y $*$ una operación binaria en $G$. El siguiente resultado establece que $G$ puede tener a lo sumo un elemento neutro. En otras palabras: \textit{el elemento neutro de un conjunto respecto a una operación binaria es único}:} \\

\begin{itemize}
    \item \Large{\textbf{Proposición:} Sean $G$ un conjunto, $e_1,e_2 \in G$ y $*$ una operación binaria en $G$. Si para todo $g \in G$, $e_1*g=g*e_1=g$ y $e_2*g=g*e_2=g$, entonces $e_1=e_2$.} \\
\end{itemize}

\Large{\textbf{Demostración:} Como $e_1 \in G$, resulta que $e_1*e_2=e_1$ (de acuerdo con la hipótesis). Similarmente: como $e_2 \in G$, resulta que $e_1*e_2=e_2$ (de acuerdo con la hipótesis). De modo pues que $e_1=e_1*e_2=e_2$, y por tanto $e_1=e_2$.} \\

\begin{flushright}
    $\square$
\end{flushright}

\Large{Como el elemento neutro respecto a una operación binaria es único, nos referimos a él de una forma especial:} \\

\begin{itemize}
    \item \Large{\textbf{Definición:} Sean $G$ un conjunto y $*$ una operación binaria en $G$. Adicionalmente, supongamos que $G$ tiene elemento neutro respecto a $*$. Nos referimos a dicho elemento como $\underline{e_G}$.} \\
\end{itemize}

\Large{\textbf{Nota:} La definición anterior solo tiene sentido si no hay ambigüedad sobre la operación binaria respecto a la cual $e_G$ es elemento neutro de $G$.} \\

\begin{flushright}
    $\blacksquare$
\end{flushright}

\begin{itemize}
    \item \Large{\textbf{Definición:} Sean $G$ un conjunto, $a,b \in G$ y $*$ una operación binaria en $G$. Adicionalmente, supongamos que $G$ tiene elemento neutro $e_G$ con respecto a $*$. Diremos que $b$ es \textbf{inverso} de $a$ con respecto a $*$, si y sólo si, $a*b=b*a=e_G$.} \\
\end{itemize}

\begin{flushright}
    $\blacksquare$
\end{flushright}

\Large{Sean $G$ un conjunto, $g \in G$ y $*$ una operación binaria en $G$. Supongamos que $G$ tiene elemento neutro $e_G$ con respecto a $*$. En general, \textit{los inversos respecto a $*$ no son únicos}. Sin embargo, si $*$ es asociativa en $G$, sí podemos asegurar la unicidad de dichos elementos:} \\

\begin{itemize}
    \item \Large{\textbf{Proposición:} Sean $G$ un conjunto, $a,b_1,b_2 \in G$ y $*$ una operación binaria en $G$. Si $G$ tiene elemento neutro $e_G$ con respecto a $*$, $*$ es asociativa en $G$ y tanto $b_1$ como $b_2$ son inversos de $a$, entonces $b_1=b_2$.} \\
\end{itemize}

\Large{\textbf{Demostración:} Note que:}
\begin{align*}
    b_1&=b_1*e_G \\
    &=b_1*(a *b_2) \\
    &=(b_1*a)*b_2 \\
    &=e_G*b_2 \\
    &=b_2
\end{align*}

\begin{flushright}
    $\square$
\end{flushright}

\Large{Como el inverso de un elemento respecto a una operación binaria asociativa es único, nos referimos a él de una forma especial:} \\

\begin{itemize}
    \item \Large{\textbf{Definición:} Sean $G$ un conjunto, $g \in G$ y $*$ una operación binaria en $G$. Adicionalmente, supongamos que $G$ tiene elemento neutro $e_G$ con respecto a $*$, $*$ es asociativa en $G$ y $g$ tiene inverso respecto a $*$. Nos referimos a dicho elemento como \underline{$g^{-1}$}.} \\
\end{itemize}

\Large{\textbf{Nota:} La definición anterior solo tiene sentido si no hay ambigüedad sobre la operación binaria respecto a la cual $g^{-1}$ es inverso de $g$.} \\

\begin{flushright}
    $\blacksquare$
\end{flushright}

\Large{Hemos introducido la noción de operación binaria y algunas de las propiedades más relevantes que puede satisfacer. Cuando una operación reúne las condiciones de asociatividad, existencia de un elemento neutro y existencia de inversos, se obtiene una estructura algebraica de especial importancia: un \textbf{grupo}. A pesar de la sencillez de sus axiomas, los grupos exhiben una enorme riqueza matemática y constituyen el objeto central de estas notas.} \\

\begin{itemize}
    \item \Large{\textbf{Definición:} Sea $G$ un conjunto y $*$ una operación binaria en $G$. Diremos que $(G,*)$ es un \textbf{grupo}, si y sólo si: \\

    i) $*$ es asociativa en $G$. \\

    ii) $G$ tiene elemento neutro respecto a $*$. \\

    \textbf{Observación:} Esta condición implica que $G \neq \emptyset$. \\

    iii) Para todo $g \in G$, $g$ tiene inverso respecto a $*$.} \\
\end{itemize}

\begin{flushright}
    $\blacksquare$
\end{flushright}

\Large{Sea $(G,*)$ un grupo. En adelante, dado que la notación ``$(G,*)$'' se refiere a la operación binaria respecto a la cual estamos trabajando, nos referimos a $e_G$ simplemente como el \textbf{elemento neutro} del grupo. Similarmente: para todo $g \in G$, nos referimos a $g^{-1}$ simplemente como el \textbf{inverso} de $g$.} \\

\Large{El siguiente teorema nos facilita un poco la tarea de mostrar que una pareja $(G,*)$ dada es un grupo:} \\

\begin{itemize}
    \item \Large{\textbf{Teorema:} Sean $G$ un conjunto y $*$ una operación binaria en $G$. Si: \\

    -- $*$ es asociativa en $G$. \\

    -- Existe $e \in G$ tal que para todo $g \in G$, $e*g=g$. \\

    -- Para todo $g \in G$, existe $g' \in G$ tal que $g' *g=e$. \\

    entonces $(G,*)$ es un grupo.} \\
\end{itemize}

\Large{\textbf{Demostración:} Sea $g \in G$. De acuerdo con la hipótesis, existe $g_1$ tal que $g_1 *g=e$. A continuación, note que:}

\begin{align*}
    g*g_1&=g*(e*g_1) \\
    &=g*[(g_1*g)*g_1] \\
    &=g*[g_1*(g*g_1)] \\
    &=(g*g_1)*(g*g_1) 
\end{align*}
\Large{Pero, nuevamente por la hipótesis, existe $g_2 \in G$ tal que $g_2*(g*g_1)=e$. Tenemos entonces que:}
\begin{align*}
    e&=g_2*(g*g_1) \\
    &=g_2*[(g*g_1)*(g*g_1)] \\
    &=[g_2*(g*g_1)]*(g*g_1) \\
    &=e*(g*g_1) \\
    &=g*g_1
\end{align*}
\Large{De modo pues que $g_1*g=g*g_1=e$. Luego:}
\begin{align*}
    g*e&=g*(g_1*g) \\
    &=(g*g_1)*g \\
    &=e*g \\
    &=g
\end{align*}
\Large{Por lo tanto, $e$ es el elemento neutro de $G$ respecto a $*$. Además, vimos que para todo $g \in G$, existe $g' \in G$ tal que $g*g'=g'*g=e$, así que cada elemento de $G$ tiene inverso respecto a $*$. Así pues, concluimos que $(G,*)$ es un grupo.} \\

\begin{flushright}
    $\square$
\end{flushright}

\Large{Sea $(G,*)$ un grupo. Cuando $*$ es conmutativa en $G$, nos referimos a $(G,*)$ con un nombre especial:} \\

\begin{itemize}
    \item \Large{\textbf{Definición:} Sea $(G,*)$ un grupo. Diremos que $(G,*)$ es \textbf{abeliano}, si y sólo si, $*$ es conmutativa en $G$.} \\
\end{itemize}

\begin{flushright}
    $\blacksquare$
\end{flushright}

\Large{\textbf{Ejemplos:}} \\

\begin{itemize}
    \item \Large{$(\mathbb Z,+)$ (donde $+$ es la suma usual) es un grupo abeliano.} \\
    
    \item{\Large{$(\mathbb Z_n,+)$ (donde $+$ es la suma usual) es un grupo abeliano.}} \\
    
    \item \Large{$(\mathbb Z_p^{\times},\cdot)$ (donde $p$ es primo y $\cdot$ es el producto usual) es un grupo abeliano.} \\
\end{itemize}

\begin{flushright}
$\blacksquare$
\end{flushright}

\Large{La cantidad de elementos de un grupo influye profundamente en sus propiedades. Por ello, resulta conveniente distinguir desde el principio entre \textbf{grupos finitos} y \textbf{grupos infinitos}:} \\

\begin{itemize}
    \item \Large{\textbf{Definición:} Sea $(G,*)$ un grupo. \\

    i) Diremos que $(G,*)$ es un \textbf{grupo finito}, si y sólo si, $G$ es finito. Además, en este caso decimos que $|G|$ es el \textbf{orden} del grupo. \\

    ii) Diremos que $(G,*)$ es un \textbf{grupo infinito}, si y sólo si, $G$ es infinito.} \\ 
\end{itemize}

\begin{flushright}
    $\blacksquare$
\end{flushright}

\Large{\textbf{Ejemplos:}} \\

\begin{itemize}
    \item \Large{($\mathbb Z,+)$ es un grupo infinito.} \\
    
    \item \Large{$(\mathbb Z_n,+)$ es un grupo finito de orden $n$.} \\
    
    \item \Large{$(\mathbb Z_p^{\times},\cdot)$ es un grupo finito de orden $p-1$.} \\
\end{itemize}

\begin{flushright}
$\blacksquare$
\end{flushright}

\end{document}